\chapter[Robust and numerical inference]{Robust estimates, information bias, and numerical stability}
\label{ch:robust-numerical}
Analysis pipelines need both statistical robustness to unusual
data points and numerical robustness to finite precision.
Neither follows merely from a mathematical formula being valid.

\section{Median and median absolute deviation}
The sample median has finite-sample replacement breakdown
point approaching $1/2$: fewer than about half the observations
cannot drive it arbitrarily far by contamination.
For observations $x_1,\ldots,x_T$, define
$\operatorname{MAD}=\operatorname{median}_i|x_i-
\operatorname{median}_j x_j|$.
For normal errors a common robust scale estimate is
$1.4826\operatorname{MAD}$, asymptotically calibrated
under the Gaussian model. For artifacted spike waveforms,
median/MAD thresholding can be more stable than a mean/standard
deviation, but physiological artifacts and nonstationarity
can contaminate even robust estimates. An influence function
formalizes the estimator's infinitesimal response to a
point-mass contamination.

\section{Empirical distribution bounds and time-rescaling}
For iid observations with distribution function $F$, let $F_T$
be the empirical CDF. A Dvoretzky--Kiefer--Wolfowitz bound
states for $\epsilon>0$
\[
\Prob\left(\sup_x|F_T(x)-F(x)|>\epsilon\right)
\le2e^{-2T\epsilon^2}.
\]
It implies uniform consistency (a Glivenko--Cantelli conclusion)
in the iid setting. The KS check of rescaled interspike intervals
uses an exponential or uniform null only when the conditional
intensity is correctly specified and relevant transformed
intervals satisfy the needed independence assumptions.
Repeated fitting and testing on the same data changes calibration.

\section{Information-theoretic chain rules and finite-sample bias}
For discrete variables, the entropy chain rule is
$H(X,Y)=H(X)+H(Y\mid X)$ and
$I(X;Y\mid Z)=H(X\mid Z)-H(X\mid Y,Z)$.
The Gaussian distribution maximizes differential entropy among
random vectors with a specified positive-definite covariance,
with entropy $\frac12\log[(2\pi e)^d\det\Sigma]$.
A plug-in entropy estimator from a finite discrete histogram
usually has downward bias. With $K$ effectively occupied bins,
a leading Miller--Madow-style correction in \emph{nats} is
$(K-1)/(2T)$ under its asymptotic sampling assumptions.
It does not cure curse-of-dimensionality or guarantee an unbiased
information estimator. Continuous estimators such as KSG have
different nearest-neighbor assumptions and biases; the
\path{mutual_information.m} implementation
should not be described as using KSG without verifying code.

\section{Conditioning, least squares, and stable algorithms}
If a full-rank matrix $X$ has spectral condition number
$\kappa_2(X)=\sigma_{\max}/\sigma_{\min}$, then
$\kappa_2(X^\top X)=\kappa_2(X)^2$.
\begin{proof}
The singular values of $X$ are the square roots of the
eigenvalues of $X^\top X$; hence those eigenvalues are
$\sigma_j^2$ and the ratio of extremes is squared.
\end{proof}
This explains why explicitly solving normal equations can
amplify roundoff relative to QR/SVD solutions. A backward-stable
algorithm returns an exact answer for a nearby problem; it does
not guarantee small forward error for ill-conditioned inputs.
The FFT reduces the computation of a DFT from the naive
$O(T^2)$ complex arithmetic to $O(T\log T)$ for appropriate
factorizable sizes. Algorithmic efficiency does not change
spectral-estimator uncertainty.

\begin{example}[Research reproducibility]\sloppy
When regressing spike history on population counts, report
condition numbers, convergence criteria, effective independent
units, model selection, bin duration and handling of omitted
variables, not just a coefficient array.
\end{example}